
    %%%%%%%%%%%%%%%%%%%%%%%%%%%%%%%%%%%%%%%%%%%%%%%%%%%%%%%%%%%%%%%%%%%%%%%%%%%%%%%%
    %%% Classe do documento
    \documentclass[12pt, addpoints]{exam}
    \UseRawInputEncoding

    %%%%%%%%%%%%%%%%%%%%%%%%%%%%%%%%%%%%%%%%%%%%%%%%%%%%%%%%%%%%%%%%%%%%%%%%%%%%%%%%
    % preambulo.tex
    %
    % Arquivo de configuração de packages para uso o LaTeX, conforme minhas
    % preferências e modelos pessoais.
    %
    % Para maiores informações, visite:
    %    
    %
    % NÃO ALTERE SE NÃO SOUBER O QUE ESTÁ FAZENDO!
    %%%%%%%%%%%%%%%%%%%%%%%%%%%%%%%%%%%%%%%%%%%%%%%%%%%%%%%%%%%%%%%%%%%%%%%%%%%%%%%%


    %%%%%%%%%%%%%%%%%%%%%%%%%%%%%%%%%%%%%%%%%%%%%%%%%%%%%%%%%%%%%%%%%%%%%%%%%%%%%%%%
    %%% Estruturas de controle:
    \input{static/utils/controles}

    %%%%%%%%%%%%%%%%%%%%%%%%%%%%%%%%%%%%%%%%%%%%%%%%%%%%%%%%%%%%%%%%%%%%%%%%%%%%%%%%
    %%% Compilação condicional em PDF:
    \input{static/utils/pdf}

    %%%%%%%%%%%%%%%%%%%%%%%%%%%%%%%%%%%%%%%%%%%%%%%%%%%%%%%%%%%%%%%%%%%%%%%%%%%%%%%%
    %%% Configurações de layout da página:
    \input{static/utils/layout}

    %%%%%%%%%%%%%%%%%%%%%%%%%%%%%%%%%%%%%%%%%%%%%%%%%%%%%%%%%%%%%%%%%%%%%%%%%%%%%%%%
    %%% Configurações para autores:
    \input{static/utils/autores}

    %%%%%%%%%%%%%%%%%%%%%%%%%%%%%%%%%%%%%%%%%%%%%%%%%%%%%%%%%%%%%%%%%%%%%%%%%%%%%%%%
    %%% Cabeçalho e rodapé:
    %%% Atenção! É MUITO DIFÍCIL ajustar apropriadamente os running headers
    %%% e running footers da primeira vez. Você pode confiar no LaTeX e deixar que
    %%% ele faça o ajuste automaticamente ou pode quebrar a cabeça e
    %%% tentar melhorar algo que o LaTeX já faz muito bem, mesmo que não fique
    %%% exatamente do jeito que você quer. O que você prefere? Se quiser ajustar
    %%% manualmente, descomente a linha abaixo e faça as configurações no arquivo
    %%% "utils/headings.tex".
    %\input{static/utils/headings}

    %%%%%%%%%%%%%%%%%%%%%%%%%%%%%%%%%%%%%%%%%%%%%%%%%%%%%%%%%%%%%%%%%%%%%%%%%%%%%%%%
    %%% Configuração para as fontes do documento:
    %%% Se você quiser usar fontes próprias ou do sistema, precisará ajustar as
    %%% configurações no arquivo "utils/fontes.tex".
    %%% ATENÇÃO: por padrão eu utilizo XeLaTeX com fontes proprietárias projetadas
    %%% por Matthew Butterick, adquiridas em https://mbtype.com, que não podem ser
    %%% distribuídas devido à licença de utilização. Se você usar XeLaTeX, você
    %%% DEVERÁ OBRIGATORIAMENTE ajustar as fontes no arquivo "utils/fontes.tex".
    \input{static/utils/fontes}

    %%%%%%%%%%%%%%%%%%%%%%%%%%%%%%%%%%%%%%%%%%%%%%%%%%%%%%%%%%%%%%%%%%%%%%%%%%%%%%%%
    %%% Sumário:
    \input{static/utils/toc}

    %%%%%%%%%%%%%%%%%%%%%%%%%%%%%%%%%%%%%%%%%%%%%%%%%%%%%%%%%%%%%%%%%%%%%%%%%%%%%%%%
    %%% Matemática:
    \input{static/utils/matematica}

    %%%%%%%%%%%%%%%%%%%%%%%%%%%%%%%%%%%%%%%%%%%%%%%%%%%%%%%%%%%%%%%%%%%%%%%%%%%%%%%%
    %%% Notas de rodapé:
    \input{static/utils/footnotes}

    %%%%%%%%%%%%%%%%%%%%%%%%%%%%%%%%%%%%%%%%%%%%%%%%%%%%%%%%%%%%%%%%%%%%%%%%%%%%%%%%
    %%% Referências cruzadas, links, citações:
    \input{static/utils/crossrefs}

    %%%%%%%%%%%%%%%%%%%%%%%%%%%%%%%%%%%%%%%%%%%%%%%%%%%%%%%%%%%%%%%%%%%%%%%%%%%%%%%%
    %%% Configurações para os hiperlinks em PDF:
    \input{static/utils/pdfconfig}

    %%%%%%%%%%%%%%%%%%%%%%%%%%%%%%%%%%%%%%%%%%%%%%%%%%%%%%%%%%%%%%%%%%%%%%%%%%%%%%%%
    %%% Glossário e índice remissivo:
    \input{static/utils/glossindex}

    %%%%%%%%%%%%%%%%%%%%%%%%%%%%%%%%%%%%%%%%%%%%%%%%%%%%%%%%%%%%%%%%%%%%%%%%%%%%%%%%
    %%% Referências bibliográficas:
    \input{static/utils/bibliografia}

    %%%%%%%%%%%%%%%%%%%%%%%%%%%%%%%%%%%%%%%%%%%%%%%%%%%%%%%%%%%%%%%%%%%%%%%%%%%%%%%%
    %%% Cores:
    \input{static/utils/cores}

    %%%%%%%%%%%%%%%%%%%%%%%%%%%%%%%%%%%%%%%%%%%%%%%%%%%%%%%%%%%%%%%%%%%%%%%%%%%%%%%%
    %%% Computação:
    \input{static/utils/computacao}

    %%%%%%%%%%%%%%%%%%%%%%%%%%%%%%%%%%%%%%%%%%%%%%%%%%%%%%%%%%%%%%%%%%%%%%%%%%%%%%%%
    %%% Gráficos:
    \input{static/utils/graficos}

    %%%%%%%%%%%%%%%%%%%%%%%%%%%%%%%%%%%%%%%%%%%%%%%%%%%%%%%%%%%%%%%%%%%%%%%%%%%%%%%%
    %%% Ícones:
    \input{static/utils/icones}

    %%%%%%%%%%%%%%%%%%%%%%%%%%%%%%%%%%%%%%%%%%%%%%%%%%%%%%%%%%%%%%%%%%%%%%%%%%%%%%%%
    %%% Blocos:
    \input{static/utils/blocos}

    %%%%%%%%%%%%%%%%%%%%%%%%%%%%%%%%%%%%%%%%%%%%%%%%%%%%%%%%%%%%%%%%%%%%%%%%%%%%%%%%
    %%% Boxes coloridos:
    \input{static/utils/colorbox}

    %%%%%%%%%%%%%%%%%%%%%%%%%%%%%%%%%%%%%%%%%%%%%%%%%%%%%%%%%%%%%%%%%%%%%%%%%%%%%%%%
    %%% Tabelas:
    \input{static/utils/tabelas}

    %%%%%%%%%%%%%%%%%%%%%%%%%%%%%%%%%%%%%%%%%%%%%%%%%%%%%%%%%%%%%%%%%%%%%%%%%%%%%%%%
    %%% Ambientes de listas:
    \input{static/utils/listas}

    %%%%%%%%%%%%%%%%%%%%%%%%%%%%%%%%%%%%%%%%%%%%%%%%%%%%%%%%%%%%%%%%%%%%%%%%%%%%%%%%
    %%% Ambientes floats e similares:
    \input{static/utils/floats}

    %%%%%%%%%%%%%%%%%%%%%%%%%%%%%%%%%%%%%%%%%%%%%%%%%%%%%%%%%%%%%%%%%%%%%%%%%%%%%%%%
    %%% Ajustes para a classe EXAM:
    \input{static/utils/exam}

    %%%%%%%%%%%%%%%%%%%%%%%%%%%%%%%%%%%%%%%%%%%%%%%%%%%%%%%%%%%%%%%%%%%%%%%%%%%%%%%%
    %%% Ajustes para textos:
    \input{static/utils/texto}

    %%%%%%%%%%%%%%%%%%%%%%%%%%%%%%%%%%%%%%%%%%%%%%%%%%%%%%%%%%%%%%%%%%%%%%%%%%%%%%%%
    %%% Ambientes, aliases, comandos e customizações em geral para serem
    %%% utilizadas nos documentos. Defina aqui qualquer customização específica
    %%% para seus documentos!
    \input{static/utils/customizacoes}

    %%%%%%%%%%%%%%%%%%%%%%%%%%%%%%%%%%%%%%%%%%%%%%%%%%%%%%%%%%%%%%%%%%%%%%%%%%%%%%%%
    %%% Ajuste do layout, espaçamento de linhas e etc.:
    \geometry{a4paper, portrait, left=2.0cm, right=2.0cm, top=2.0cm, bottom=2.0cm}
    \singlespacing

    %%%%%%%%%%%%%%%%%%%%%%%%%%%%%%%%%%%%%%%%%%%%%%%%%%%%%%%%%%%%%%%%%%%%%%%%%%%%%%%%
    %%% Configurações de cabeçalho e rodapé (não use fancyhdr pois ocorre conflito):
    \pagestyle{headandfoot}
    \runningheadrule
    \firstpageheader{}{}{}
    \runningheader{Sem disciplina}{}{Avaliação}
    \firstpagefooter{}{}{Página \thepage\ de \numpages}
    \runningfooter{}{}{Página \thepage\ de \numpages}
    \runningfootrule

    %%%%%%%%%%%%%%%%%%%%%%%%%%%%%%%%%%%%%%%%%%%%%%%%%%%%%%%%%%%%%%%%%%%%%%%%%%%%%%%%
    %%% Configurações para as propriedades do PDF:
    \hypersetup{
    hidelinks,           % Comente para web, descomente para impressão
    colorlinks=true,      % True para web, False para impressão
    pdftitle=Prova Pre Enade: Avaliação,
    pdfauthor=Professor,
    pdfsubject=Sem disciplina,
    pdfkeywords=['Sem tag'],
    pdfinfo={
        CreationDate={}, % Ex.: D:AAAAMMDDHH24MISS
        ModDate={}       % Ex.: D:AAAAMMDDHH24MISS
    }
    }
    \input{static/utils/pdfconfig}

    %%%%%%%%%%%%%%%%%%%%%%%%%%%%%%%%%%%%%%%%%%%%%%%%%%%%%%%%%%%%%%%%%%%%%%%%%%%%%%%%
    %%% Começa o documento
    \begin{document}

    %%%%%%%%%%%%%%%%%%%%%%%%%%%%%%%%%%%%%%%%%%%%%%%%%%%%%%%%%%%%%%%%%%%%%%%%%%%%%%%%
    %%% Folha de instruções padronizadas da UVV para a prova
    \pdfbookmark[1]{Instruções}{instrucoes}
    \begin{figure}[H]
        \begin{center}
            \includegraphics[scale=0.3]{{static/imgs/logo-uvv.png}}
        \end{center}	
    \end{figure}

    \vspace{-1cm}

    \begin{table}[H]
        \centering
        \begin{tabular}{|llll|}
            \hline
            \multicolumn{2}{|l|}{\textbf{Disciplina:} Sem disciplina} & 
            \multicolumn{1}{l|}{\multirow{3}{*}{\textbf{\makecell{Nota:\\ \,\\ \,}}}} & 
            \multirow{3}{*}{\textbf{\makecell{Coordenador:\\ \,\\ \,}}} \\ 
            \cline{1-2}
            \multicolumn{2}{|l|}{\textbf{Professor:} Professor \hspace{4.72cm} } & 
            \multicolumn{1}{l|}{} & \\ 
            \cline{1-2}
            \multicolumn{2}{|l|}{\textbf{Aluno:}} & 
            \multicolumn{1}{l|}{} & \\ 
            \hline
            \multicolumn{1}{|l|}{\textbf{Turma:} \hspace{6.3cm} } & 
            \multicolumn{1}{l|}{\textbf{Semestre:}} & 
            \multicolumn{2}{l|}{\textbf{Valor:} 10 pontos} \\ 
            \hline
            \multicolumn{1}{|l|}{\textbf{Data:}} & 
            \multicolumn{3}{l|}{\textbf{Avaliação:}} \\ 
            \hline
        \end{tabular}
    \end{table}

    \begin{center}
        \fbox{\setlength\fboxsep{0.3cm} \fbox{\parbox{15.4cm}{
            \begin{center}
                \textbf{INSTRUÇÕES PARA A REALIZAÇÃO DESTA AVALIAÇÃO:}
            \end{center}
            \begin{itemize}[leftmargin=*]
                \item Esta é a primeira avaliação da disciplina Sem disciplina, 
                    e engloba o conteúdo referente Sem tag.
                \item Esta avaliação contém 40 questões objetivas, \textbf{todas obrigatórias}.
                \item \textbf{Leia atentamente cada questão!} As questões objetivas terão uma,
                    e apenas uma única, resposta a ser marcada.
                \item \textbf{Desligue o celular} antes de começar. A avaliação é
                    \textbf{individual} e \textbf{sem consulta}.
                \item As questões podem ser respondidas, na prova, com lápis ou caneta. Ao final
                    da prova \textbf{{VOCÊ DEVERÁ PREENCHER O GABARITO COM CANETA
                    DE COR PRETA}} \textbf{\underline{OU AZUL}} para que seja feita a correção
                    automatizada através da leitura do padrão de respostas marcado no
                    gabarito.
                \item \textbf{CUIDADO ao preencher o gabarito} pois eles são \textbf{INDIVIDUAIS
                    e NOMEADOS} (cada estudante tem um gabarito próprio específico, com um
                    código de identificação único). \textbf{Questões rasuradas são anuladas}
                    pelo software de leitura do gabarito.
                \item Ao final da prova, \textbf{DEVOLVA A PROVA E O GABARITO} para o professor.
                \item Siga todas as normas de \textbf{Integridade Acadêmica} da disciplina.
                    Alunos que forem flagrados com qualquer espécie de ``cola'' ou trocando
                    informações com outros alunos terão suas avaliações recolhidas, as notas
                    zeradas, e a situação será encaminhada para a coordenação para a aplicação
                    das penalidades previstas pela universidade.
                \item Com 90 minutos de avaliação você terá tempo de até 2,min por questão,
                    em média.
                \item Boa avaliação!
            \end{itemize}
            \vspace{2cm}
        }}}
    \end{center}
    \newpage

    %%%%%%%%%%%%%%%%%%%%%%%%%%%%%%%%%%%%%%%%%%%%%%%%%%%%%%%%%%%%%%%%%%%%%%%%%%%%%%%%
    %%% Inicia o bloco de questões:
    \begin{questions}

    \newpage
    \question

Sistemas de processamento de transações, tais como sistemas de reservas aéreas, devem prover um mecanismo que garanta que cada transação não é afetada por outras transações que possam estar ocorrendo ao mesmo tempo. Transações de duas fases obedecem a um protocolo que garante essa atomicidade. Em transações de duas fases:


\begin{choices}
\choice Todas as operações de leitura ocorrem antes da primeira operação de escrita.
\choice Qualquer objeto corretamente travado deve ser destravado antes que outro objeto possa ser travado.
\choice Todas as ações de travamento (\textit{lock}) ocorrem antes da primeira ação de destravamento.
\choice Verifica-se a disponibilidade de todas as travas antes de executar qualquer ação de travamento.
\choice Uma trava compartilhada sobre um objeto deve ser obtida antes de uma trava exclusiva sobre o objeto ser obtida.
\end{choices}

\question

A derivada da função \( f(x) = x^x \) é igual a:
\begin{choices}
\choice \( x x^{x-1} \)
\choice \( x^x \ln(x) \)
\choice \( x^x (\ln(x) + x) \)
\choice \( x^x \)
\choice \( x^x (\ln(x) + 1) \)
\end{choices}

\question

Professor Mac Sperto propôs o seguinte algoritmo de ordenação, chamado de Super Merge, similar ao \textit{merge sort}: divida o vetor em 4 partes do mesmo tamanho (ao invés de 2, como é feito no \textit{merge sort}). Ordene recursivamente cada uma das partes e depois intercale-as por um procedimento semelhante ao procedimento de intercalação do \textit{merge sort}. Qual das alternativas abaixo é verdadeira?
\begin{choices}
\choice Super Merge está correto, mas consome tempo \( O(\text{merge sort}) \)
\choice Super Merge está correto, mas consome tempo maior que \( O(\text{merge sort}) \)
\choice Super Merge não está correto. Não é possível ordenar quebrando o vetor em 4 partes
\choice Super Merge está correto, mas consome tempo menor que \( O(\text{merge sort}) \)
\choice Nenhuma das afirmações acima está correta
\end{choices}

\question

São motivos importantes e gerais para o estudo das linguagens de programação,

exceto:
\begin{choices}
\choice Melhorar o embasamento na escolha das linguagens para um projeto
\choice Entender a importância da implementação
\choice Aumentar a capacidade de expressar idéias
\choice Avanço geral da computação
\choice Aprender a sintaxe e a semântica de uma determinada linguagem
\end{choices}

\question

O usuário que não é da área de tecnologia, costuma ser de alta gerência, que é

especialista no negócio e tem a responsabilidade central pelas políticas de

dados da empresa é o:
\begin{choices}
\choice Gerente de Tecnologia da Informação (ITM: IT Manager)
\choice Administrador do Banco de Dados (DBA: Database Administrator)
\choice Encarregado da Proteção de Dados (DPO: Data Protection Officer)
\choice Administrador de Dados (DA: Data Administrator)
\choice Coordenador de Tecnologia da Informação (ITS: IT Supervisor)
\end{choices}

\question

Em relação ao teste de software, qual das afirmações a seguir é INCORRETA: 
\begin{choices}
\choice Um teste bem sucedido é aquele que revela um erro ainda não descoberto. 
\choice Um bom caso de teste é aquele que tem uma elevada probabilidade de revelar um erro ainda não descoberto. 
\choice O processo de depuração é a parte mais imprevisível do processo de teste pois um erro pode demorar uma hora, um dia ou um mês para ser diagnosticado e corrigido. 
\choice Os dados compilados quando a atividade de teste é levada a efeito proporcionam uma boa indicação da confiabilidade do software e alguma indicação da qualidade do software como um todo. 
\choice A atividade de teste é o processo de executar um programa com a intenção de demonstrar a ausência de erros. 
\end{choices}

\question

Considere as seguintes afirmações sobre resolução de problemas em IA.

\begin{enumerate}[label=\Roman*.]

    \item A* é um conhecido algoritmo de busca heurística.

    \item O Minimax é um dos principais algoritmos para jogos de dois jogadores, como o xadrez.

    \item Busca em espaço de estados é uma das formas de resolução de problemas em IA.

\end{enumerate}

São corretas:
\begin{choices}
\choice Apenas II e III
\choice Apenas I e II
\choice Apenas I e III
\choice I, II e III
\choice Apenas III
\end{choices}

\question

Um compilador detecta:
\begin{choices}
\choice erros nos resultados gerados pelo programa
\choice erros de sintaxe do programa 
\choice erros que podem ocorrer durante a execução do programa
\choice erros aritméticos
\choice todos os erros citados acima
\end{choices}

\question

Redes semânticas, frames e lógica são formalismos utilizados principalmente em: 
\begin{choices}
\choice inferência em sistemas especialistas
\choice redes neurais
\choice IA distribuída
\choice descoberta de conhecimento em bases de dados
\choice representação de conhecimento
\end{choices}

\question

``Um modelo de dados é uma definição lógica, abstrata e auto-contida

dos \circled{1}, \circled{2} e \circled{3} que, juntos, constituem a \circled{4}

com a qual os usuários interagem. Os \circled{1} nos permitem modelar

a \circled{5}; os \circled{2} nos permitem modelar

seu \circled{6}. Os \circled{3} nos permitem modelar sua \circled{7}.'



As palavras/termos que completam a frase acima, são, conforme os números:


\begin{choices}
\choice 1 = demais conceitos; 2 = objetos; 3 = operadores; 4 = estrutura de dados; 5 = máquina abstrata; 6 = integridade; 7 = comportamento.
\choice 1 = objetos; 2 = operadores; 3 = máquina abstrata; 4 = demais conceitos; 5 = estrutura de dados; 6 = integridade; 7 = comportamento.
\choice 1 = operadores; 2 = objetos; 3 = demais conceitos; 4 = máquina abstrata; 5 = estrutura de dados; 6 = comportamento; 7 = integridade.
\choice 1 = objetos; 2 = máquina abstrata; 3 = demais conceitos; 4 = operadores; 5 = estrutura de dados; 6 = comportamento; 7 = integridade.
\choice 1 = objetos; 2 = operadores; 3 = demais conceitos; 4 = máquina abstrata; 5 = estrutura de dados; 6 = comportamento; 7 = integridade.
\end{choices}

\question

No que diz respeito as vantagens da arquitetura de micro-núcleo para sistemas operacionais em relação a arquiteturas de núcleo monolítico, quais das seguintes afirmações são verdadeiras?



\begin{itemize}

    \item[I] A arquitetura de micro-núcleo facilita a depuração do SO.

    \item[II] A arquitetura de micro-núcleo permite um número menor de mudanças de contexto.

    \item[III] A arquitetura de micro-núcleo facilita a reconfiguração de serviços do SO pois a maioria deles reside em espaço de usuário.

\end{itemize}
\begin{choices}
\choice II e III
\choice Todas são verdadeiras
\choice Apenas I
\choice I e II
\choice I e III
\end{choices}

\question

Considere as seguintes afirmações sobre mecanismos de inferência em sistemas baseados em regras.

\begin{enumerate}[label=\Roman*.]

    \item O encadeamento regressivo tem pouca utilidade prática, pois deve partir do possível resultado.

    \item O encadeamento progressivo tanto pode ser em amplitude quanto em profundidade.

    \item Podem trabalhar com informações incertas ou incompletas. 

\end{enumerate}

São corretas:
\begin{choices}
\choice Apenas II e III
\choice Apenas I e II
\choice I, II e III
\choice Apenas III
\choice Apenas I e III
\end{choices}

\question

Sobre a técnica conhecida como Z-buffer é correto afirmar que:
\begin{choices}
\choice Nenhuma das alternativas acima está correta. 
\choice É possível realizar o cômputo das variáveis envolvidas de forma incremental.
\choice As dimensões do Z-buffer são independentes das dimensões do frame buffer. 
\choice É uma técnica muito comum de detecção de colisão. 
\choice As primitivas geométricas precisam estar ordenadas de acordo com a distância em relação ao observador. 
\end{choices}

\question

Considere as seguintes afirmativas sobre as linguagens usadas para análise sintática:

\begin{enumerate}[label=\Roman*.]

    \item A classe LL(1) não aceita linguagens com produções que apresentem recursões diretas à esquerda (ex. \(L \rightarrow La\)) mas aceita linguagens com recursões indiretas (ex. \(L \rightarrow Ra, R \rightarrow Lb\)).

    \item A linguagem LR(1) reconhece a mesma classe de linguagens que LALR(1).

    \item A linguagem SLR(1) reconhece uma classe de linguagens maior que LR(0).

\end{enumerate}

Selecione a afirmativa correta:
\begin{choices}
\choice As afirmativas I e III são verdadeiras
\choice As afirmativas I e III são falsas
\choice As afirmativas II e III são verdadeiras
\choice As afirmativas I e II são verdadeiras
\choice Apenas a afirmativa III é verdadeira
\end{choices}

\question

Quanto ao TCP, é INCORRETO afirmar:
\begin{choices}
\choice Possui um campo de \textit{checksum} que valida as informações de seu cabeçalho, mas não valida as informações de \textit{payload} (campo de dados).
\choice É um protocolo do nível de transporte.
\choice Utiliza portas para permitir a comunicação entre processos localizados em dispositivos diferentes.
\choice É um protocolo orientado a conexão.
\choice Usa janelas deslizantes para implementar o controle de fluxo e erro.
\end{choices}

\question

O pipeline de visualização de objetos tridimensionais reúne um conjunto de transformações e processos aplicados a primitivas geométricas. Sobre essas transformações e processos pode-se dizer que:

\begin{enumerate}[label=\Roman*.]

    \item Os objetos devem corresponder a sólidos.

    \item As coordenadas dos vértices sofrem transformação de acordo com a posição e orientação do observador.

    \item Um volume de visualização correspondente a um paralelepípedo é determinado pela adoção de projeção perspectiva.

    \item A fase final do pipeline corresponde à rasterização dos polígonos.

\end{enumerate}

Selecione a alternativa correta:
\begin{choices}
\choice Apenas a afirmativa IV é falsa.
\choice Apenas as afirmativas I e III são falsas.
\choice As afirmativas II e III são falsas.
\choice Todas as afirmativas são verdadeiras.
\choice Apenas a afirmativa IV está verdadeira.
\end{choices}

\question

Qual das seguintes afirmações sobre crescimento assintótico de funções não é verdadeira:
\begin{choices}
\choice Se \( f(n) = O(g(n)) \) então \( g(n) = O(f(n)) \)
\choice \( 2n^2 + 3n + 1 = O(n^2) \)
\choice Se \( f(n) = O(g(n)) \) e \( g(n) = O(h(n)) \) então \( f(n) = O(h(n)) \)
\choice \( 2^{n+1} = O(2^n) \)
\choice \( \log n^2 = O(\log n) \)
\end{choices}

\question

Em relação aos metadados de um banco de dados, assinale a afirmação correta:
\begin{choices}
\choice São armazenados no catálogo (ou dicionário de dados) do sistema de gerenciamento de bancos de dados, que fica nas mesmas tabelas de dados dos usuários, permitindo assim que os metadados estejam sempre atualizados.
\choice Armazenam informações como os dados de cadastro dos usuários, os valores de compra em uma ordem de compra, as turmas e disciplinas que os professores dão aula, e dados de recursos humanos.
\choice Como são sempre atualizados de forma automática pelo sistema de gerenciamento de bancos de dados, estão sempre prontos para nos dar uma ``foto' da estrutura do banco de dados em um momento no tempo.
\choice São atualizados pelo SGBD, sempre que ocorre alguma operação de inserção, atualização ou remoção de dados.
\choice Armazenam informações sobre os dados dos usuários, ou seja, são dados de rotina dos usuários, aqueles dados que os usuários precisam para trabalhar no seu dia a dia.
\end{choices}

\question

Sobre a hierarquia de Chomsky podemos afirmar que:
\begin{choices}
\choice Uma linguagem que é recursivamente enumerável não pode ser uma linguagem regular
\choice As linguagens reconhecidas por autômatos a pilha são as linguagens regulares
\choice As linguagens livres de contexto e as linguagens sensíveis a contexto se excluem
\choice Uma linguagem que não é regular é livre de contexto
\choice Há linguagens que não são nem livres de contexto nem sensíveis a contexto
\end{choices}

\question

Assinale a alternativa INCORRETA:


\begin{choices}
\choice O controle de fluxo tem como objetivo garantir que nenhum dos parceiros de uma comunicação inunda o outro enviando pacotes mais rápido do que ele pode tratar.
\choice Serviços orientados a conexão podem ser implementados em subredes que funcionam no modo datagrama.
\choice Nos serviços orientados a conexões há a necessidade de estabelecimento de uma conexão antes da transferência dos dados.
\choice Os serviços orientados a conexão podem ajudar no controle de congestionamento através da diminuição da taxa de transmissão durante um congestionamento em andamento.
\choice Os serviços orientados a conexões são sempre confiáveis garantindo a entrega ordenada e completa dos dados transmitidos.
\end{choices}

\question

Sejam os seguintes predicados de uma linguagem de primeira ordem:



\begin{itemize}

    \item \( N(x) \): \( x \) é número;

    \item \( P(x) \): \( x \) tem propriedade \( P \);

    \item \( x < y \): \( x \) é menor que \( y \).

\end{itemize}



E sejam os símbolos:

\begin{itemize}

    \item \( \forall \): quantificador universal;

    \item \( \Rightarrow \): operador se-então;

    \item \( \neg \): operador de negação.

\end{itemize}



Para a fórmula: \( \forall x (N(x) \Rightarrow \neg \forall y (N(y) \Rightarrow y < x)) \), qual alternativa abaixo \textbf{NÃO} constitui uma tradução possível?


\begin{choices}
\choice Para todo número, existe um outro número que é maior do que ele.
\choice Não há um número tal que todos os números são menores do que ele.
\choice Para todo número, não é verdade que qualquer número seja menor do que ele.
\choice Não há um número menor do que outro número.
\choice Para qualquer \( x \), se \( x \) é número, então não é verdade que todos os números são menores do que ele.
\end{choices}

\question

Para a gramática a seguir, qual o conjunto de terminais que pode aparecer como primeiro terminal após o não-terminal \( A \), em qualquer forma sentencial gerada pela gramática abaixo (isto é, não necessariamente imediatamente após \( A \)), onde \( \epsilon \) representa a sentença vazia?



$S \rightarrow ABCDd$\\

$A \rightarrow aA \mid \epsilon$\\

$B \rightarrow bC \mid \epsilon$\\

$C \rightarrow cD \mid \epsilon$\\

$D \rightarrow e$\\
\begin{choices}
\choice \{b, c, d, e\}
\choice \{b\}
\choice \{b, c, e\}
\choice \{e\}
\choice \{d\}
\end{choices}

\question

Considere as seguintes tabelas em uma base de dados relacional (chaves primárias sublinhadas):\\

Departamento (\underline{CodDepto}, NomeDepto)\\

Empregado (\underline{CodEmp}, NomeEmp, CodDepto)\\

Considere as seguintes restrições de integridade sobre esta base de dados relacional:\\

- Empregado.CodDepto é sempre diferente de NULL\\

- Empregado.CodDepto é chave estrangeira da tabela Departamento com cláusulas ON DELETE RESTRICT e ON UPDATE RESTRICT\\

Qual das seguintes validações não é especificada por estas restrições de integridade:
\begin{choices}
\choice Sempre que uma nova linha for inserida em Empregado, deve ser garantido que o valor de Empregado.CodDepto aparece na coluna Departamento.CodDepto.
\choice Sempre que o valor de Departamento.CodDepto for alterado, deve ser garantido que não há uma linha com o antigo valor de Departamento.CodDepto na coluna Empregado.CodDepto.
\choice Sempre que o valor de Empregado.CodDepto for alterado, deve ser garantido que o novo valor de Empregado.CodDepto aparece em Departamento.CodDepto.
\choice Sempre que uma linha for excluída de Departamento, deve ser garantido que o valor de Departamento.CodDepto não aparece na coluna Empregado.CodDepto.
\choice Sempre que uma nova linha for inserida em Departamento, deve ser garantido que o valor de Departamento.CodDepto aparece na coluna Empregado.CodDepto.
\end{choices}

\question

Dentre as características do modelo relacional e do modelo de objetos em bancos de dados, qual afirmação é \textbf{INCORRETA}?


\begin{choices}
\choice O relacionamento de herança é diretamente representado no modelo relacional.
\choice O relacionamento binário \( N \times M \) é representado de modo semelhante nos dois modelos.	
\choice O modelo de objetos provê uma representação bem próxima de linguagens de programação.
\choice O modelo de objetos é mais adequado para a representação de tipos abstratos de dados.
\choice O modelo de objetos possui mais recursos estruturais para a representação de dados que o relacional.
\end{choices}

\question

Qual o significado de coerência de memórias \textit{cache} em sistemas multiprocessados?


\begin{choices}
\choice Caches em processadores diferentes nunca compartilham a mesma linha de cache.
\choice Caches em processadores diferentes sempre contêm o mesmo dado válido para a mesma linha de cache.
\choice Caches em processadores diferentes sempre lêem os mesmos dados ao mesmo tempo.
\choice Caches em processadores diferentes podem possuir dados diferentes associados à mesma linha de cache.
\choice Caches em processadores diferentes nunca interagem entre si.
\end{choices}

\question

Na Álgebra Booleana:
\begin{choices}
\choice Os dígitos são octais, de 0 a 7
\choice Há dez valores motivados pelos dez dedos do ser humano
\choice Os dígitos são hexadecimais de 0 a 15
\choice Os dígitos são alfanuméricos que podem ser representados por um byte
\choice Os dígitos são binários 0 e 1
\end{choices}

\question

O conjunto básico de atividades e a ordem em que são realizadas no processo de construção de um software definem o que é habitualmente denominado de ciclo de vida do software. O ciclo de vida tradicional (também denominado waterfall) ainda é hoje em dia um dos mais difundidos e tem por característica principal: 
\begin{choices}
\choice a codificação de uma versão executável do sistema desde as fases iniciais do desenvolvimento, de modo que o sistema final é incrementalmente construído, daí a alusão à ideia de "cascata" (waterfall); 
\choice o uso de formalização rigorosa em todas as etapas de desenvolvimento; 
\choice a priorização da análise dos riscos do desenvolvimento
\choice a avaliação constante dos resultados intermediários feita pelo cliente.
\choice a abordagem sistemática para realização das atividades do desenvolvimento de software de modo que elas seguem um fluxo sequencial; 
\end{choices}

\question Assinale a alternativa \textbf{incorreta}:
\begin{choices}
\choice Um SGBD não tem interface gráfica
\choice No modelo relacional, tudo o que o usuário percebe são tabelas
\choice Modelo de dados é a mesma coisa que o diagrama relacional
\choice O modelo relacional de dados foi criado por Edgar Frank Cood
\choice O MariaDB é um SGBD
\end{choices}

\question

A interposição de um circuito de memória cache entre o processador e a memória principal (RAM)
\begin{choices}
\choice aumenta o tráfego de instruções e/ou dados entre memória e disco
\choice permite acessos concorrentes à memória RAM
\choice diminui o tráfego de instruções e/ou dados no barramento de memória
\choice aumenta o tráfego de instruções e/ou dados no barramento de memória
\choice diminui o tráfego de instruções e/ou dados entre memória e disco
\end{choices}

\question

Analise o diagrama relacional do banco de dados exibido na figura abaixo:

\begin{figure}[H]

  \centering

  \caption{Banco de dados de peças e fornecedores}

  \label{fig:pecas}

  %\fbox{

    \includegraphics[width=0.5\linewidth]{static/imgs/defaul.png}

  %}\\

  %\footnotesize{Fonte: xxxx.}

\end{figure}

Podemos dizer que esse banco de dados ilustra:
\begin{choices}
\choice Um relacionamento com cardinalidade N:N entre as tabelas ``pecas' e ``fornecimentos', realizado através de relacionamentos não identificados com a tabela ``fornecedores'. 
\choice Um relacionamento com cardinalidade N:N entre as tabelas ``pecas' e ``fornecedores', realizado através de relacionamentos não identificados através da tabela ``fornecimentos'. 
\choice Um relacionamento com identificação N:N entre as tabelas ``pecas' e ``fornecedores', realizado através do relacionamento com cardinalidade identificada através da tabela ``fornecimentos'. 
\choice Um relacionamento com cardinalidade N:N entre as tabelas ``pecas' e ``fornecedores', realizado através de relacionamentos identificados através da tabela ``fornecimentos'. 
\choice Um relacionamento com cardinalidade 1:N entre as tabelas ``pecas' e ``fornecedores', indicando que uma peça pode ter sido fornecida por diversos fornecedores. 
\end{choices}

\question

Considere as seguintes afirmações sobre redes neurais artificiais: \begin{enumerate}[label=\Roman*.] \item Um perceptron elementar só computa funções linearmente separáveis. \item Não aceitam valores numéricos como entrada. \item O "conhecimento" é representado principalmente através do peso das conexões. \end{enumerate} São corretas:
\begin{choices}
\choice Apenas I e III
\choice Apenas III
\choice Apenas II e III
\choice Apenas I e II 
\choice I, II e III
\end{choices}

\question

Ao analisar um novo sistema de gerenciamento de bancos de dados que surgiu no

mercado, você percebe que a entrada às operações do modelo de dados são tabelas

mas as saídas nunca são tabelas, são sempre outras estruturas. O vendedor desse

SGBD afirma que ele é um SGBD relacional. Pode-se afirmar que:


\begin{choices}
\choice Esse SGBD não é relacional pois, para ser verdadeiramente relacional, as operações do modelo de dados não podem trabalhar recebendo tabelas como entrada.
\choice Esse SGBD é relacional pois consegue trabalhar com tabelas para a entrada das operações do modelo de dados. A saída não é importante.
\choice Esse SGBD é relacional pois o vendedor afirmou que é.
\choice Nenhuma das alternativas anteriores.
\choice Esse SGBD não é relacional pois no modelo relacional a única estrutura percebida pelos usuários são tabelas.
\end{choices}

\question

O processo de visualização de objetos 3D envolve uma série de passos desde a representação vetorial de um objeto até a exibição da imagem correspondente na tela do computador (pipeline 3D). Selecione a alternativa abaixo que reflete a ordem correta em que esses passos devem ocorrer. 
\begin{choices}
\choice Transformação de câmera, mapeamento para coordenadas de tela, recorte 3D, rasterização, projeção. 
\choice Recorte 3D, transformação de câmera, rasterização, projeção, mapeamento para coordenadas de tela. 
\choice Transformação de câmera, recorte 3D, projeção, mapeamento para coordenadas de tela, rasterização. 
\choice Projeção, transformação de câmera, recorte 3D, mapeamento para coordenadas de tela, rasterização. 
\choice Nenhuma das respostas acima está correta.
\end{choices}

\question

Seja a seguinte linguagem, onde \(\epsilon\) representa o string vazio e \$ representa um marcador de fim de entrada:



$S \rightarrow ABCD \\$\\

$A \rightarrow a \mid \epsilon \\$\\

$B \rightarrow a \mid \epsilon \\$\\

$C \rightarrow c \mid \epsilon \\$\\

$D \rightarrow S \mid c \mid \epsilon$\\



É incorreto afirmar que:
\begin{choices}
\choice O conjunto FIRST(A) = a, \(\epsilon\)
\choice O conjunto FOLLOW(B) = c, \$
\choice O conjunto FOLLOW(A) = a, c, \$
\choice O conjunto FIRST(D) é igual ao conjunto FIRST(S)
\choice O conjunto FOLLOW(D) é igual a FOLLOW(S)
\end{choices}

\question

Suponha que \( T \) seja uma árvore AVL inicialmente vazia, e considere a inserção dos elementos \(10, 20, 30, 5, 15, 2\) em \( T \), nesta ordem. Qual das sequências abaixo corresponde a um percurso de \( T \) em pré-ordem:
\begin{choices}
\choice 20, 10, 5, 2, 15, 30
\choice 15, 10, 5, 2, 20, 30
\choice 10, 5, 2, 20, 15, 30
\choice 2, 5, 10, 15, 20, 30
\choice 30, 20, 15, 10, 5, 2
\end{choices}

\question

\begin{figure}[H]

    \centering

    \includegraphics[width=0.5\linewidth]{static/imgs/defaul.png}

    \caption{Enter Caption}

\end{figure}

Se \( O = (0, 0, 0) \); \( A = (2, 4, 1) \); \( B = (3, 1, 1) \) e \( C = (1, 3, 5) \) então o volume do sólido acima é


\begin{choices}
\choice 30
\choice 35
\choice 44
\choice 35/2
\choice 21
\end{choices}

\question

Marque a alternativa onde todos os conceitos estão corretos.


\begin{choices}
\choice Em um diagrama de fluxo de dados, uma entidade externa representa um produtor ou um consumidor de informação e está fora dos limites do sistema modelado; cada processo pode ser refinado, para explicitar um maior detalhamento; um DFD contém dois níveis de detalhamento; um processo é um transformador de informação e também está fora do sistema; o nível 0 de um DFD representa o sistema como um todo e indica os principais usuários e as funções do sistema.
\choice Em um diagrama de fluxo de dados uma entidade externa representa uma fonte ou destino das informações processadas pelo sistema e está fora dos limites do sistema modelado; cada processo pode ser refinado, para explicitar um maior detalhamento; um DFD pode conter vários níveis de detalhamento; um processo é um transformador de informação; o nível 0 de um DFD representa o sistema como um todo e indica as principais fontes e destinos das informações, usualmente referenciado por Diagrama de Contexto.
\choice Nenhuma das alternativas anteriores.
\choice Em um diagrama de fluxo de dados uma entidade externa representa um produtor ou um consumidor de informação e está fora dos limites do sistema modelado; cada processo deve ser refinado, para explicitar um maior detalhamento; um DFD pode conter vários níveis de detalhamento; um processo é um transformador de informação e também está fora do sistema; o nível 0 de um DFD representa o sistema como um todo e indica os principais usuários e as funções do sistema.
\choice Em um diagrama de fluxo de dados uma entidade externa representa uma fonte ou destino das informações processadas pelo sistema e está fora dos limites do sistema modelado; cada processo pode ser refinado, para explicitar um maior detalhamento; um DFD pode conter vários níveis de detalhamento; um processo é um transformador de informação e também está fora do sistema; o nível 0 de um DFD representa o sistema como um todo e indica as principais fontes e destinos das informações.
\end{choices}

\question

Em uma lista circular duplamente encadeada com \( n \) elementos, o espaço ocupado apenas pelos apontadores é (assuma que um apontador ocupa \( p \) bytes):
\begin{choices}
\choice \( 4np \)
\choice \( 2np \)
\choice \( (np)^2 \)
\choice \( 6np \)
\choice \( np \)
\end{choices}

\question

São vantagens da utilização de threads no espaço do usuário, exceto:
\begin{choices}
\choice Maior portabilidade da aplicação.
\choice Nenhuma modificação é necessária no kernel.
\choice A criação e o gerenciamento das threads são mais eficientes.
\choice O sistema operacional escalona a thread.
\choice O escalonamento pode ser específico para a aplicação.
\end{choices}

\question

Algoritmos distribuídos podem usar passagem de "token" por um anel lógico para implementar exclusão mútua ou ordenação global de mensagens. Nesses algoritmos, apenas o processo que possui o "token" tem a permissão de usar um recurso compartilhado ou numerar mensagens, por exemplo. Considerando o conceito acima, podemos afirmar que: 
\begin{choices}
\choice para usar essa abordagem, os computadores precisam estar conectados em uma rede com topologia em 
\choice a abordagem deve tratar no mínimo dois tipos de defeitos: perda do "token" e colapso de processos 
\choice a abordagem é adequada apenas para sistemas onde possa ser controlado o tempo que cada computador permanece com o 
\choice nessa abordagem é impossível evitar a geração espontânea de vários "tokens" mesmo em sistemas livres de falhas 
\choice a abordagem é pouco robusta pois a perda do "token" por um processo provoca o bloqueio do algoritmo distribuído que a usa 
\end{choices}



    %%%%%%%%%%%%%%%%%%%%%%%%%%%%%%%%%%%%%%%%%%%%%%%%%%%%%%%%%%%%%%%%%%%%%%%%%%%%%%%%
    %%% Finaliza o bloco de questões:
    \end{questions}
    \end{document}
    